% 초록. 본문을 다 쓴 뒤 마지막에 손본다.

\begin{abstract}
본 보고서는 서울시 토지거래허가구역(이하 토허제) 지정이 아파트 \emph{거래량과
가격}에 미친 인과효과를 실거래 미시자료로 추정한다. 국토교통부 아파트 매매
실거래 \emph{상세} 자료 47만 건(분석 8개 자치구, 2006-01$\sim$2026-05)과 서울시보·
전자관보의 지정 공고 원문을 결합해 \textbf{거래 한 건마다 계약일·전용면적·층·단지
식별자와 그 시점의 규제 상태를 붙인 미시 패널}을 구축하고, 2020년 6월 지정,
2021년 4월 확대 지정, 2025년 2월 부분 해제와 3월 재지정을 사건으로 삼아
이중차분·사건연구·삼중차분·합성대조군으로 추정했다.

주요 발견은 넷이다.
(1) \textbf{지정 직후 3개월 거래량이 $-71 \sim -80\%$다.} 법정동 27개 가운데 4개를
처치로 고르는 모든 경우 17,550가지 중 진짜 배정이 1위이고, 사전 추세 위배를
실제 관측값의 1.6배까지 열어 두어도 구간이 0을 넘지 않으며, 합리적인 설계 조합
192가지가 전부 음수다. \textbf{열 달 뒤 완전히 다른 네 자치구에서 일어난
2021년 4월 지정에서도 같은 크기였고}, 두 사건의 차이는 0과 구분되지 않는다
($+8.7\%$, $p=0.715$).
(2) \textbf{선행연구끼리 결론이 엇갈린 이유는 지역 정의가 아니라 사건 시점
정렬이었다.} 설계 축별로 분해하면 사건 시점 정렬이 39.6\,\%p, 통제군 정의가
2.3\,\%p를 가른다. 달력 월로 자르면 $-40.5\%$, 효력일에 맞추면 $-80.1\%$다.
(3) \textbf{효과는 모두에게 같지 않았고, 그 차이는 대출 규제가 아니라 허가 대상의
면적 기준에서 왔다.} 2020년 지정은 대지 18\,m$^2$ 초과만 허가 대상이어서 소형은
거래가 줄지 않았고(전용 40\,m$^2$ 미만 $+125.7\%$), 2022년 6월 기준이 6\,m$^2$로
강화되자 \emph{그때} 소형만 33.2\% 줄었다. 기준을 바꾸지 않은 2021년 재지정에서는
소형이 움직이지 않았다($-8.1\%$, $p=0.61$). 이 메커니즘에서 세운
\textbf{사전 예측이 2025년 재지정 자료에서 확인됐다}(소형 $-71.5\%$).
(4) \textbf{가격을 잡았다는 증거는 약하다.} 평행추세를 가정하면 $-4 \sim -7\%$
(헤도닉+단지 $-5.9\%$, 반복매매 $-6.0\%$)이나, 그 가정을 조금만 풀면 0과
구분되지 않고(붕괴값 0.22) 합성대조군으로는 확인되지 않는다($p=0.511$).

\textbf{말할 수 없는 것도 적는다.} 효과가 얼마나 오래 가는지는 기준 분기 선택에
따라 $-20 \sim +16\%$로 흔들려 이 설계로 답할 수 없다. 부과와 해제의 대칭성도
가릴 수 없다. 일반화 범위는 두 사건·아홉 개 동이며 모두 서울 고가 재건축
지역이므로 \emph{서울 고가 재건축 지역에서는 일관된다}까지가 정확하다.

본 보고서의 또 하나의 축은 \textbf{분석의 전제가 되는 시스템}이다. 원자료 수집부터
그림·표 생성까지 \pipeSteps 단계 파이프라인 하나로 돌고, 문서 수치 검증·고아
산출물 점검·의존 점검·표 생성 넷이 매 실행마다 걸린다. 새 폴더에 내려받아 새 가상환경에서
돌려 결과 JSON 25개가 바이트 단위로 같은 것을 확인했다. 코드와 결정기록 \pipeDecisions 편
전체를 공개한다: \url{https://github.com/cosmicpotato2047/capstone-ltz-causal}.
\end{abstract}

\vspace{1em}
\noindent\textbf{키워드:} 토지거래허가구역, 이중차분, 사건연구, 미시 패널,
부동산 정책, 재현가능 연구

\vspace{1em}
\begin{center}\textbf{Abstract (English)}\end{center}
\small
This report estimates the causal effect of Seoul's Land Transaction Permit
Zones (LTZ) on apartment \emph{transaction volume and prices} using
transaction-level administrative data. We combine 472,782 detailed apartment
sale records (8 districts, 2006-01 to 2026-05) with the original designation
notices published in the Seoul and national gazettes, constructing a
\textbf{micro panel in which every transaction carries its contract date, floor
area, floor level, complex identifier, and the regulatory status in force on
that day}. Using the June 2020 designation, the April 2021 expansion, and the
February 2025 partial release followed by the March 2025 re-designation as
events, we estimate difference-in-differences, event-study, triple-difference,
and synthetic-control specifications.

Four findings.
(1) \textbf{Transaction volume falls by 71--80\% in the three months after
designation.} Among all 17{,}550 ways of drawing 4 treated neighbourhoods from
27, the true assignment ranks first; the result survives opening the
pre-trend violation to 1.6 times its observed magnitude
\citep{rambachan2023honest}; and all 192 reasonable specification combinations
are negative. \textbf{The April 2021 designation, ten months later and in four
entirely different districts, produced an effect of the same size}, and the
difference between the two events is indistinguishable from zero
($+8.7\%$, $p=0.715$).
(2) \textbf{Prior studies disagree because of event-time alignment, not
geography.} Decomposing by design axis, event-time alignment accounts for
39.6\,pp of the spread while the control-group definition accounts for 2.3\,pp.
Aligning on calendar months gives $-40.5\%$; aligning on the effective date
gives $-80.1\%$.
(3) \textbf{The effect was not uniform, and the heterogeneity came from the
permit-threshold floor area rather than from mortgage regulation.} Under the
2020 rule only lots above 18\,m$^2$ required a permit, so small units were
exempt and their transactions did not fall (units under 40\,m$^2$: $+125.7\%$).
When the threshold tightened to 6\,m$^2$ in June 2022, small units fell by
33.2\% \emph{at that moment}, while the 2021 re-designation, which left the
threshold unchanged, produced no movement ($-8.1\%$, $p=0.61$).
A \textbf{prediction registered from this mechanism was confirmed} in the 2025
re-designation (small units $-71.5\%$).
(4) \textbf{Evidence that prices were contained is weak.} Assuming parallel
trends gives $-4$ to $-7\%$, but relaxing that assumption slightly makes the
interval include zero (breakdown value 0.22) and the synthetic control does not
confirm it ($p=0.511$).

We also state what cannot be said: the duration of the effect swings between
$-20\%$ and $+16\%$ depending on the reference period, and the symmetry between
imposition and release cannot be resolved. External validity extends to two
events and nine neighbourhoods, all high-priced redevelopment areas in Seoul.

A second axis of this report is the \textbf{system that makes the analysis
possible}: a single \pipeSteps-step pipeline from data collection to figures and
tables, with automated checks for stale numbers in documents, orphaned
artefacts, dependency drift, and table generation. A clean-room reproduction in a fresh clone and
virtual environment produced 25 result files identical byte for byte. Code and
all \pipeDecisions{} decision records are public.
\normalsize
